\subsection{Multiplicative Structure on K(C)}

Let K be a commutative ring and A a bigebra over the ring K (III, §11, No.~4, p.~585), with coproduct c and counit \(\gamma\) . Unless stated otherwise, the tensor products are over K. Let E and F be (left) A-modules. The tensor product \(E \otimes F\) is endowed with an \((A \otimes A)\) -module structure characterized by the formula

\[
(a \otimes b) (x \otimes y) = a x \otimes b y\tag{11}
\]

for \(a, b \in \mathrm{A}\) , \(x \in \mathrm{E}\) , and \(y \in \mathrm{F}\) . Using the homomorphism \(c: \mathrm{A} \to \mathrm{A} \otimes \mathrm{A}\) , we deduce from this \((\mathrm{A} \otimes \mathrm{A})\) -module an A-module \(c_*(\mathrm{E} \otimes \mathrm{F})\) (II, §1, No.~13, p.~221). More precisely, let \(a \in \mathrm{A}\) ; if \(c(a) = \sum_{i} a_{i} \otimes b_{i}\) , then we have

\[
a (x \otimes y) = \sum_ {i} a _ {i} x \otimes b _ {i} y\tag{12}
\]

for \(x \in E\) and \(y \in F\) . By abuse of notation, we also denote the resulting A-module by \(E \otimes F\) . It follows immediately from the coassociativity of c that the canonical isomorphism of K-modules

\[
\varphi : (\mathrm{E} \otimes \mathrm{F}) \otimes \mathrm{G} \longrightarrow \mathrm{E} \otimes (\mathrm{F} \otimes \mathrm{G})
\]

is A-linear for all A-modules E, F, and G. Likewise, if the bigebra A is cocommutative, then the canonical isomorphism from \(E \otimes F\) to \(F \otimes E\) is A-linear. Finally, let \(K_{\gamma}\) be the A-module with underlying group K and external law \((a, x) \mapsto \gamma(a)x\) . The canonical isomorphism from \(K \otimes E\) (resp. \(E \otimes K\) ) to E is an isomorphism of A-modules from \(K_{\gamma} \otimes E\) (resp. \(E \otimes K_{\gamma}\) ) to E.

PROPOSITION 9. --- Let K be a commutative ring, A a bigebra over K with counit γ, and C an additive set of classes of A-modules with the following properties:

\begin{enumerate}
\def\labelenumi{(\roman{enumi})}
\item
  Every A-module of type \(\mathcal{C}\) is a projective (* or, more generally, flat*) K-module.
\item
  If the A-modules E and F are of type C, then so is the A-module \(E \otimes F\) .
\item
  The A-module \(\mathrm{K}_{\gamma}\) defined above is of type \(\mathcal{C}\) .
\end{enumerate}

Then there exists a unique ring structure on the additive group \(\mathrm{K}(\mathcal{C})\) whose multiplication satisfies

\[
[ \mathrm{E} ] _ {\mathcal {C}} [ \mathrm{F} ] _ {\mathcal {C}} = [ \mathrm{E} \otimes \mathrm{F} ] _ {\mathcal {C}}
\]

for all A-modules E and F of type C. The unit element of \(\mathrm{K}(\mathcal{C})\) is \([\mathrm{K}_{\gamma}]_{\mathcal{C}}\) . If the bigebra A is cocommutative, then the ring \(\mathrm{K}(\mathcal{C})\) is commutative.

Endowed with the law of composition given by \((\mathrm{E},\mathrm{F})\mapsto\mathrm{cl}(\mathrm{E}\otimes\mathrm{F})\) , the set C is a monoid with unit element \(\mathrm{cl}(\mathrm{K}_{\gamma})\) . Consequently, \(Z^{(\mathcal{C})}\) is canonically endowed with the structure of a ring with multiplication characterized by the formula

\[
\operatorname{cl} (\mathrm{E}) \operatorname{cl} (\mathrm{F}) = \operatorname{cl} (\mathrm{E} \otimes \mathrm{F})\tag{13}
\]

and unit element \(\operatorname{cl}(\mathrm{K}_{\gamma})\) (III, §2, No.~6, p.~446).

Given an A-module F of type C and an exact sequence

\[
0 \longrightarrow \mathrm{E} ^ {\prime} \longrightarrow \mathrm{E} \longrightarrow \mathrm{E} ^ {\prime \prime} \longrightarrow 0\tag{\( \mathcal{E} \)}
\]

of A-modules of type C, the sequence

\[
0 \longrightarrow \mathrm{E} ^ {\prime} \otimes \mathrm{F} \longrightarrow \mathrm{E} \otimes \mathrm{F} \longrightarrow \mathrm{E} ^ {\prime \prime} \otimes \mathrm{F} \longrightarrow 0\tag{\(\left(\mathcal{E}\otimes \mathrm{F}\right)\)}
\]

deduced from \((\mathcal{E})\) is exact because F is projective ( \(^{*}\) or, more generally, flat \(_{*}\) ) over K (II, §3, No.~6, p.~251, Proposition 5 and II, §3, No.~7, p.~257, Corollary 6). In the notation of No.~2, we then have

\[
r _ {\mathcal {E} \otimes \mathrm{F}} = r _ {\mathcal {E}} \operatorname{cl} (\mathrm{F}).\tag{14}
\]

It follows that the subgroup R of \(\mathbf{Z}^{(\mathcal{C})}\) generated by the elements of the form \(r_{\mathcal{E}}\) is a right ideal of the ring \(\mathbf{Z}^{(\mathcal{C})}\) , and one proves likewise that it is a left ideal of \(\mathbf{Z}^{(\mathcal{C})}\) . By definition, \(\mathrm{K}(\mathcal{C})\) is the quotient group \(\mathbf{Z}^{(\mathcal{C})}/\mathrm{R}\) ; hence there exists a unique ring structure on \(\mathrm{K}(\mathcal{C})\) whose multiplication satisfies \([E]_{\mathcal{C}}[F]_{\mathcal{C}} = [E \otimes F]_{\mathcal{C}}\) for all A-modules E and F of type C. Its unit element is \([K_{\gamma}]\) .

When the bigebra A is cocommutative, the monoid C is commutative; it follows that the ring \(\mathbf{Z}^{(\mathcal{C})}\) and the quotient ring \(\mathrm{K}(\mathcal{C})\) are commutative.

Remark. --- Under only assumptions (i) and (ii) of Proposition 9, the Grothendieck group \(\mathrm{K}^{\prime}(\mathcal{C})\) for direct sums (VIII, p.~188) has a unique ring structure whose multiplication satisfies \([E]_{\mathcal{C}}^{\prime}[F]_{\mathcal{C}}^{\prime} = [E \otimes F]_{\mathcal{C}}^{\prime}\) . Its unit element is \([K_{\gamma}]_{\mathcal{C}}^{\prime}\) . The ring \(\mathrm{K}^{\prime}(\mathcal{C})\) is commutative if the bigebra A is cocommutative. The proof is analogous to that of Proposition 9 because the group \(\mathrm{K}^{\prime}(\mathcal{C})\) can be identified with \(\mathbf{Z}^{(\mathcal{C})}/\mathrm{R}^{\prime}\) , where \(R^{\prime}\) is the subgroup of \(\mathbf{Z}^{(\mathcal{C})}\) generated by the elements of the form \(\mathrm{cl}(\mathrm{E}^{\prime} \oplus \mathrm{E}^{\prime\prime}) - \mathrm{cl}(\mathrm{E}^{\prime}) - \mathrm{cl}(\mathrm{E}^{\prime\prime})\) (loc. cit.).

Under assumptions (i), (ii), and (iii) of Proposition 9, the ring \(\mathrm{K}(\mathcal{C})\) is called the Grothendieck ring of C. These conditions are, in particular, satisfied when K is a field and C is the set of classes of A-modules that are finite-dimensional over K. Consequently, we have the following.

COROLLARY. --- Let A be a bigebra with counit \(\gamma\) over a commutative field K. Then there exists a unique ring structure on the additive group \(\mathrm{R}_{\mathrm{K}}(\mathrm{A})\) whose multiplication satisfies

\[
[ \mathrm{E} ] _ {\mathcal {C}} [ \mathrm{F} ] _ {\mathcal {C}} = [ \mathrm{E} \otimes_ {\mathrm{K}} \mathrm{F} ] _ {\mathcal {C}}
\]

for all A-modules E and F that are finite-dimensional over K. The unit element of \(\mathrm{R}_{\mathrm{K}}(\mathrm{A})\) is \([\mathrm{K}_{\gamma}]_{\mathcal{C}}\) . If the bigebra A is cocommutative, then the ring \(\mathrm{R}_{\mathrm{K}}(\mathrm{A})\) is commutative.

Examples. --- 1) Let K be a commutative field. Let G be a group, and let K{[}G{]} be the algebra of the group G. We identify G with its canonical image in K{[}G{]} (III, §2, No.~6, p.~446).

We endow K{[}G{]} with the structure of a bigebra with coproduct c and counit \(\gamma\) given by

\[
c (g) = g \otimes g, \quad \gamma (g) = 1 \quad (g \in G).\tag{15}
\]

Let \(\mathrm{E}\) and \(\mathrm{F}\) be \(\mathrm{K}[\mathrm{G}]\) -modules; by formula (12), the \(\mathrm{K}[\mathrm{G}]\) -module structure on \(\mathrm{E} \otimes \mathrm{F}\) is given by

\[
g (x \otimes y) = g x \otimes g y \quad (g \in \mathrm{G}, x \in \mathrm{E}, y \in \mathrm{F}).\tag{16}
\]

The K{[}G{]}-module \(K_{\gamma}\) is the vector space K endowed with the action of G defined by \(g\lambda = \lambda\) for \(g \in G\) and \(\lambda \in K\) .

The ring \(\mathrm{R}_{\mathrm{K}}(\mathrm{K}[\mathrm{G}])\) is also denoted by \(\mathrm{R}_{\mathrm{K}}(\mathrm{G})\) . It is commutative; its multiplication is given by {[}E{]} {[}F{]} = {[}E \(\otimes_{\mathrm{K}}\) F{]}, and the unit element of \(\mathrm{R}_{\mathrm{K}}(\mathrm{G})\) is \([K_{\gamma}]\) .

*2) Let \(\mathfrak{g}\) be a Lie algebra over a commutative field \(K\) and \(U(\mathfrak{g})\) its enveloping algebra; we identify \(\mathfrak{g}\) with its canonical image in \(U(\mathfrak{g})\) (Lie, I, §2, No.~7, p.~39, Corollary 2). We endow \(U(\mathfrak{g})\) with the structure of a bigebra with coproduct \(c\) and counit \(\gamma\) given by

\[
c (\xi) = \xi \otimes 1 + 1 \otimes \xi , \quad \gamma (\xi) = 0\tag{17}
\]

for \(\xi\in\mathfrak{g}\) (Lie, II, §1, No.~4, p.~115).

Let \(\mathrm{E}\) and \(\mathrm{F}\) be \(\mathrm{U}(\mathfrak{g})\) -modules; by formula (17), the \(\mathrm{U}(\mathfrak{g})\) -module structure on \(\mathrm{E} \otimes \mathrm{F}\) is characterized by

\[
\xi (x \otimes y) = \xi x \otimes y + x \otimes \xi y\tag{18}
\]

for \(\xi\in\mathfrak{g}, x\in\mathrm{E}\) , and \(y\in\mathrm{F}\) .

The Grothendieck ring \(\mathrm{R}_{\mathrm{K}}(\mathrm{U}(\mathfrak{g}))\) is also denoted by \(\mathcal{R}(\mathfrak{g})\) in Lie, VIII, §7, No.~6, p.~139.*

Let A be a commutative ring. We can view A as a bigebra that is cocommutative over itself, with coproduct the natural isomorphism from A to \(A \otimes_{A} A\) and counit \(Id_{A}\) (III, §11, No.~4, p.~585). By Proposition 9, we obtain the following result.

PROPOSITION 10. --- Let A be a commutative ring, and let C be an additive set of classes of A-modules satisfying the following three conditions:

\begin{enumerate}
\def\labelenumi{(\roman{enumi})}
\item
  Every A-module of type \(\mathcal{C}\) is projective (* or, more generally, flat*)
\item
  If \(\mathrm{E}\) and \(\mathrm{F}\) are A-modules of type \(\mathcal{C}\) , then the A-module \(\mathrm{E} \otimes_{\mathrm{A}} \mathrm{F}\) is also of type \(\mathcal{C}\) .
\item
  The A-module A is of type \(\mathcal{C}\) .
\end{enumerate}

Then there exists a unique ring structure on the additive group \(\mathrm{K}(\mathcal{C})\) satisfying \([E]_{\mathcal{C}}[F]_{\mathcal{C}} = [E \otimes_{A} F]_{\mathcal{C}}\) for every pair E, F of A-modules of type C. The unit element of \(\mathrm{K}(\mathcal{C})\) is \([A]_{\mathcal{C}}\) .

